\begin{abstract}
Vision-Language Models (VLMs) frequently err not because they reason incorrectly,
but because they reason from \emph{incorrect visual premises}---false perceptual
beliefs about what an image contains.
We introduce \textbf{Belief-Constrained Inference} (\bci{}), an inference-time
framework that externalizes, verifies, and corrects a model's intermediate visual
beliefs before final answer generation, without modifying model weights.

\bci{} prompts the VLM to emit a set of atomic, falsifiable visual claims as an
explicit reasoning substrate. These claims are verified against structured visual
evidence (scene graphs) and, if contradicted, are replaced with ground-truth facts
prior to constrained re-reasoning. On GQA ($n{=}500$), \bci{} with full premise
replacement (\textbf{E1}) raises accuracy from 59.2\% to 77.2\% ($+18.0$ pp;
McNemar $p{=}1.5{\times}10^{-14}$)---a causal upper bound establishing premise
correction as a highly effective intervention axis.
Scaled to $n{=}2000$, E1 achieves $+16.2$ pp (63.55\% to 79.75\%;
$p{=}1.78{\times}10^{-48}$) with a flip-to-correct rate of 21.2\%, confirming
consistent gains at scale.
A practical confidence-gated variant (\textbf{E12b}, $n{=}2000$) achieves
$+3.45$ pp ($p{=}1.3{\times}10^{-5}$) while intervening on only 59.85\% of
samples, with multi-seed stability of $66.18\%\pm 0.14$ pp.
Critically, targeted correction outperforms random belief deletion by $+35$ pp,
confirming that gains are claim-specific rather than a generic perturbation effect.
\bci{} also outperforms matched lightweight baselines including self-correction
(49.2\%) and chain-of-verification (57.2\%), establishing it as a strong and
interpretable inference-time reliability mechanism for VLMs.
\end{abstract}
